\documentclass[10pt,a4paper]{amsart}
\usepackage{amsmath,amssymb,amsthm,graphicx}
% Retained minimal preamble for the Stacks Project chapter flat.tex (More on Flatness).
% Verbatim excerpt of the chapter preamble: Xy-pic setup (native preamble lines 14-18) and the native
% theorem declarations with the mathematical macros (native preamble lines 154-217). Copied mechanically;
% nothing added, removed or redefined. The wrapper class and the amsmath/amssymb/amsthm/graphicx base
% packages are supplied by the assembler, so the class/externaldocument part of the native preamble is omitted.
% For commutative diagrams we use Xy-pic
\usepackage[all]{xy}

% We use 2cell for 2-commutative diagrams.
\xyoption{2cell}
% Theorem environments.
%
\theoremstyle{plain}
\newtheorem{theorem}[subsection]{Theorem}
\newtheorem{proposition}[subsection]{Proposition}
\newtheorem{lemma}[subsection]{Lemma}

\theoremstyle{definition}
\newtheorem{definition}[subsection]{Definition}
\newtheorem{example}[subsection]{Example}
\newtheorem{exercise}[subsection]{Exercise}
\newtheorem{situation}[subsection]{Situation}

\theoremstyle{remark}
\newtheorem{remark}[subsection]{Remark}
\newtheorem{remarks}[subsection]{Remarks}

\numberwithin{equation}{subsection}

% Macros
%
\def\lim{\mathop{\mathrm{lim}}\nolimits}
\def\colim{\mathop{\mathrm{colim}}\nolimits}
\def\Spec{\mathop{\mathrm{Spec}}}
\def\Hom{\mathop{\mathrm{Hom}}\nolimits}
\def\Ext{\mathop{\mathrm{Ext}}\nolimits}
\def\SheafHom{\mathop{\mathcal{H}\!\mathit{om}}\nolimits}
\def\SheafExt{\mathop{\mathcal{E}\!\mathit{xt}}\nolimits}
\def\Sch{\mathit{Sch}}
\def\Mor{\mathop{\mathrm{Mor}}\nolimits}
\def\Ob{\mathop{\mathrm{Ob}}\nolimits}
\def\Sh{\mathop{\mathit{Sh}}\nolimits}
\def\NL{\mathop{N\!L}\nolimits}
\def\CH{\mathop{\mathrm{CH}}\nolimits}
\def\proetale{{pro\text{-}\acute{e}tale}}
\def\etale{{\acute{e}tale}}
\def\QCoh{\mathit{QCoh}}
\def\Ker{\mathop{\mathrm{Ker}}}
\def\Im{\mathop{\mathrm{Im}}}
\def\Coker{\mathop{\mathrm{Coker}}}
\def\Coim{\mathop{\mathrm{Coim}}}

% Boxtimes
%
\DeclareMathSymbol{\boxtimes}{\mathbin}{AMSa}{"02}

%
% Macros for moduli stacks/spaces
%
\def\QCohstack{\mathcal{QC}\!\mathit{oh}}
\def\Cohstack{\mathcal{C}\!\mathit{oh}}
\def\Spacesstack{\mathcal{S}\!\mathit{paces}}
\def\Quotfunctor{\mathrm{Quot}}
\def\Hilbfunctor{\mathrm{Hilb}}
\def\Curvesstack{\mathcal{C}\!\mathit{urves}}
\def\Polarizedstack{\mathcal{P}\!\mathit{olarized}}
\def\Complexesstack{\mathcal{C}\!\mathit{omplexes}}
% \Pic is the operator that assigns to X its picard group, usage \Pic(X)
% \Picardstack_{X/B} denotes the Picard stack of X over B
% \Picardfunctor_{X/B} denotes the Picard functor of X over B
\def\Pic{\mathop{\mathrm{Pic}}\nolimits}
\def\Picardstack{\mathcal{P}\!\mathit{ic}}
\def\Picardfunctor{\mathrm{Pic}}
\def\Deformationcategory{\mathcal{D}\!\mathit{ef}}

% --- appended 2026-10-12 (Chief): the project own definitions of the reference/slogan/history
% display environments, copied mechanically from the shared Stacks preamble lines 8-12
% (\usepackage{verbatim} plus the three \newenvironment lines, which the chapter preamble keeps
% verbatim from the project preamble). Nothing added, removed or redefined: these environments
% carry no mathematical content and the project itself typesets them as comments. Without these
% three lines a counted statement whose own text uses them fails with Environment slogan undefined.
\usepackage{verbatim}
\newenvironment{reference}{\comment}{\endcomment}
\newenvironment{slogan}{\comment}{\endcomment}
\newenvironment{history}{\comment}{\endcomment}
\usepackage{cleveref}
\crefname{originalequation}{original equation}{original equations}
\Crefname{originalequation}{Original equation}{Original equations}
\makeatletter
\newlabel{equation-almost-blow-up-square}{{0EV6}{source}{}{sourcepointer-1}{}}
\newlabel{equation-almost-blow-up-square@cref}{{[originalequation][0][]0EV6}{[1][0][]source}{}{}{}}
\makeatother
\crefname{originallemma}{original lemma}{original lemmas}
\Crefname{originallemma}{Original lemma}{Original lemmas}
\makeatletter
\newlabel{morphisms-lemma-base-change-proper}{{01W4}{source}{}{sourcepointer-2}{}}
\newlabel{morphisms-lemma-base-change-proper@cref}{{[originallemma][0][]01W4}{[1][0][]source}{}{}{}}
\makeatother
\crefname{originallemma}{original lemma}{original lemmas}
\Crefname{originallemma}{Original lemma}{Original lemmas}
\makeatletter
\newlabel{morphisms-lemma-base-change-finite-presentation}{{01TS}{source}{}{sourcepointer-3}{}}
\newlabel{morphisms-lemma-base-change-finite-presentation@cref}{{[originallemma][0][]01TS}{[1][0][]source}{}{}{}}
\makeatother
\crefname{originallemma}{original lemma}{original lemmas}
\Crefname{originallemma}{Original lemma}{Original lemmas}
\makeatletter
\newlabel{morphisms-lemma-base-change-closed-immersion}{{01QR}{source}{}{sourcepointer-4}{}}
\newlabel{morphisms-lemma-base-change-closed-immersion@cref}{{[originallemma][0][]01QR}{[1][0][]source}{}{}{}}
\makeatother
\crefname{originallemma}{original lemma}{original lemmas}
\Crefname{originallemma}{Original lemma}{Original lemmas}
\makeatletter
\newlabel{divisors-lemma-pullback-locally-principal}{{053P}{source}{}{sourcepointer-5}{}}
\newlabel{divisors-lemma-pullback-locally-principal@cref}{{[originallemma][0][]053P}{[1][0][]source}{}{}{}}
\makeatother
\crefname{originaldefinition}{original definition}{original definitions}
\Crefname{originaldefinition}{Original definition}{Original definitions}
\makeatletter
\newlabel{divisors-definition-strict-transform}{{080D}{source}{}{sourcepointer-6}{}}
\newlabel{divisors-definition-strict-transform@cref}{{[originaldefinition][0][]080D}{[1][0][]source}{}{}{}}
\makeatother
\crefname{originallemma}{original lemma}{original lemmas}
\Crefname{originallemma}{Original lemma}{Original lemmas}
\makeatletter
\newlabel{divisors-lemma-strict-transform}{{080E}{source}{}{sourcepointer-7}{}}
\newlabel{divisors-lemma-strict-transform@cref}{{[originallemma][0][]080E}{[1][0][]source}{}{}{}}
\makeatother
\title{Almost blowup squares are stable under base change}
\author{The Stacks Project Authors}
\date{stacks-project:flat.tex:local-snapshot-20261011; GNU FDL 1.2 or later}
\begin{document}
\maketitle
\noindent\textit{Editorial note.} All delivered mathematical text is the author's own, copied from the pinned Stacks Project chapter file flat.tex; no statement or proof step was rewritten, reordered or removed.\par\medskip
\noindent\textit{Editorial note.} The preamble is the chapter's own retained minimal preamble (sources/stacks-flat/retained-preamble.tex), a verbatim excerpt of the native Xy-pic setup and of the native theorem declarations and mathematical macros, followed by the project own definitions of the reference, slogan and history display environments (verbatim package plus three newenvironment lines) that the retained file carries at its end and that the project itself typesets as comments. No macro was invented or redefined, and the wrapper class and the amsmath/amssymb/amsthm/graphicx base packages are supplied by the assembler.\par\medskip
\noindent\textit{Editorial note.} The counted statement and proof contain no cross reference, no citation and no included file, so the delivered body is self contained.\par\medskip
\noindent\textit{Editorial note.} Canonical location of the counted proposition: Stacks tag 0EV7 (sources/stacks-flat/pinned-original-source.tex line 10987).\par\medskip
\noindent\textit{Allowed standard prerequisites (editorial).}\begin{itemize}
\item The defining properties of an almost blow up square. Same chapter, equation labelled equation-almost-blow-up-square and its accompanying definition, invoked in the statement. Lists the properties verified in the proof.
\item Base change of proper morphisms, finite presentation, and closed immersions. Morphisms, Lemmas labelled morphisms-lemma-base-change-proper, morphisms-lemma-base-change-finite-presentation and morphisms-lemma-base-change-closed-immersion, cited at proof lines 10999-11005. Transfers the first two properties to the base changed square.
\item Local principality of the pullback of a locally principal closed subscheme. Divisors, Lemma labelled divisors-lemma-pullback-locally-principal, cited at proof lines 11007-11009. Gives the third property.
\item Strict transforms and their identification with blowups. Divisors, Definition labelled divisors-definition-strict-transform and Lemma labelled divisors-lemma-strict-transform, cited at proof lines 11015-11021. Identifies Y'' as the blow up of Y x\_X Z in Y.
\end{itemize}\medskip
\ifdefined\hypertarget\hypertarget{sourcepointer-1}{}\fi
\noindent\textit{Source reference.} Original equation 0EV6: sources/stacks-flat/pinned-original-source.tex, source lines 10957--10957. Same-chapter pointer into the author environment equation at source line 10957; invoked at proof line None of the counted proof. It is retained only as a navigation pointer; every load-bearing step is inside the delivered span.\par\medskip
\ifdefined\hypertarget\hypertarget{sourcepointer-2}{}\fi
\noindent\textit{Source reference.} Original lemma 01W4: sources/stacks-chow/tags.txt, source lines 2469--2469. Cross-chapter standard prerequisite located in the project tag list staged in the local snapshot; invoked at proof line 11001 of the counted proof. It is retained only as a navigation pointer; every load-bearing step is inside the delivered span.\par\medskip
\ifdefined\hypertarget\hypertarget{sourcepointer-3}{}\fi
\noindent\textit{Source reference.} Original lemma 01TS: sources/stacks-chow/tags.txt, source lines 2385--2385. Cross-chapter standard prerequisite located in the project tag list staged in the local snapshot; invoked at proof line 11002 of the counted proof. It is retained only as a navigation pointer; every load-bearing step is inside the delivered span.\par\medskip
\ifdefined\hypertarget\hypertarget{sourcepointer-4}{}\fi
\noindent\textit{Source reference.} Original lemma 01QR: sources/stacks-chow/tags.txt, source lines 2276--2276. Cross-chapter standard prerequisite located in the project tag list staged in the local snapshot; invoked at proof line 11004 of the counted proof. It is retained only as a navigation pointer; every load-bearing step is inside the delivered span.\par\medskip
\ifdefined\hypertarget\hypertarget{sourcepointer-5}{}\fi
\noindent\textit{Source reference.} Original lemma 053P: sources/stacks-chow/tags.txt, source lines 6612--6612. Cross-chapter standard prerequisite located in the project tag list staged in the local snapshot; invoked at proof line 11008 of the counted proof. It is retained only as a navigation pointer; every load-bearing step is inside the delivered span.\par\medskip
\ifdefined\hypertarget\hypertarget{sourcepointer-6}{}\fi
\noindent\textit{Source reference.} Original definition 080D: sources/stacks-chow/tags.txt, source lines 10178--10178. Cross-chapter standard prerequisite located in the project tag list staged in the local snapshot; invoked at proof line 11018 of the counted proof. It is retained only as a navigation pointer; every load-bearing step is inside the delivered span.\par\medskip
\ifdefined\hypertarget\hypertarget{sourcepointer-7}{}\fi
\noindent\textit{Source reference.} Original lemma 080E: sources/stacks-chow/tags.txt, source lines 10179--10179. Cross-chapter standard prerequisite located in the project tag list staged in the local snapshot; invoked at proof line 11019 of the counted proof. It is retained only as a navigation pointer; every load-bearing step is inside the delivered span.\par\medskip
\begin{lemma}
\label{lemma-base-change-almost-blow-up}
Consider an almost blow up square (\ref{equation-almost-blow-up-square}).
Let $Y \to X$ be any morphism. Then the base change
$$
\xymatrix{
Y \times_X E \ar[d] \ar[r] & Y \times_X X' \ar[d] \\
Y \times_X Z \ar[r] & Y
}
$$
is an almost blow up square too.
\end{lemma}

\begin{proof}
The morphism $Y \times_X X' \to Y$ is proper and of finite presentation
by Morphisms, Lemmas \ref{morphisms-lemma-base-change-proper} and
\ref{morphisms-lemma-base-change-finite-presentation}.
The morphism $Y \times_X Z \to Y$ is a closed immersion
(Morphisms, Lemma \ref{morphisms-lemma-base-change-closed-immersion}) of
finite presentation. The inverse image of $Y \times_X Z$ in $Y \times_X X'$
is equal to the inverse image of $E$ in $Y \times_X X'$ and hence is
locally principal
(Divisors, Lemma \ref{divisors-lemma-pullback-locally-principal}).
Let $X'' \subset X'$, resp.\ $Y'' \subset Y \times_X X'$ be the closed
subscheme corresponding to the quasi-coherent ideal of sections of
$\mathcal{O}_{X'}$, resp.\ $\mathcal{O}_{Y \times_Y X'}$
supported on $E$, resp.\ $Y \times_X E$.
Clearly, $Y'' \subset Y \times_X X''$ is the closed subscheme
corresponding to the quasi-coherent ideal of sections of
$\mathcal{O}_{Y \times_Y X''}$ supported on $Y \times_X (E \cap X'')$.
Thus $Y''$ is the strict transform of $Y$ relative to the blowing up
$X'' \to X$, see
Divisors, Definition \ref{divisors-definition-strict-transform}.
Thus by Divisors, Lemma \ref{divisors-lemma-strict-transform}
we see that $Y''$ is the blow up of $Y \times_X Z$ on $Y$.
\end{proof}

\end{document}
